\subsection{Virtual punctures and the characteristic-zero fibre construction}

\begin{claim}{C.LIT1}{Proved: the exact pullback construction}{C4}
Let $k$ have characteristic zero, let $C/k$ be hyperbolic, and consider
the projection $C^{(2)}=(C\times_k C)\setminus\Delta_C\to C$.
Work in a specified quotient category in which there is an exact
sequence
\[
 1\longrightarrow F\longrightarrow E\xrightarrow{q}\Pi_C
 \longrightarrow1.
\]
For every section $s:G_k\to\Pi_C$ the profinite pullback
\[
 \Pi_s=E\times_{\Pi_C,s}G_k
\]
is an extension $1\to F\to\Pi_s\to G_k\to1$. If $s=s_P$ is the
section of $P\in C(k)$, the natural fibre map identifies $\Pi_s$ with
the corresponding quotient of $\Pi_{C\setminus P}$, whenever the
assumed exact fibration statement includes that fibre identification.
\end{claim}
\begin{proof}
The kernel of the pullback projection is $F\times\{1\}$. Surjectivity
follows from that of $q$, proving exactness. For a geometric section
the map from the arithmetic group of the actual fibre has the specified
geometric fibre map and is the identity on $G_k$. If those geometric
groups agree, the extension isomorphism criterion proves that the
arithmetic groups agree. These arguments are valid for profinite groups
and continuous maps, without requiring the section to come from a point.
\end{proof}

\paragraph{The quotient qualification is necessary.}
An exact sequence in a quotient cannot recover information already
discarded by that quotient. For example, killing the entire diagonal
inertia gives the product group with its exact projection sequence;
its fibre is the unpunctured geometric group. Consequently the phrase
``the punctured group in every exact quotient'' means the corresponding
punctured quotient and its retained boundary data. It cannot literally
mean the full punctured fundamental group in a quotient that kills
the new cusp. This is a limitation of that wording, not a failure of
the pullback construction.

\begin{claim}{C.LIT2}{Proved: retaining the distinguished cusp}{C4}
Assume the chosen configuration quotient also retains a diagonal
boundary decomposition group $D_\Delta$ mapping surjectively to
$\Pi_C$, with kernel the surviving inertia $I_\Delta$. Then
\[
 D_s=D_\Delta\times_{\Pi_C,s}G_k\subset\Pi_s
\]
is the distinguished virtual cusp extension. The pair $(\Pi_s,D_s)$
depends only on the geometric conjugacy class of $s$, up to the
corresponding inner ambiguity. For $s_P$ it is the actual punctured
fibre group with its new cusp in the chosen quotient.
\end{claim}
\begin{proof}
Suppose $s'=\Inn(\gamma)s$ for $\gamma\in\Delta_C$. The diagonal
decomposition map has a geometric lift $d$ of $\gamma$; conjugation
by $d$ identifies the inverse images of the two section subgroups and
their intersections with $D_\Delta$. Thus it transports both the
extension and the cusp. If the second proper projection is retained,
it changes that projection by $\Inn(\gamma)$, which may be removed
by a geometric inner automorphism in the target fibre because its
map to the other proper factor is surjective. This leaves the cusp
conjugacy class unchanged. For $s_P$, a local parameter at $P$ gives
the ordinary boundary decomposition description, and the exact fibre
comparison gives the asserted identification.
\end{proof}

The diagonal lift is extra geometric information. Exactness of
$E\to\Pi_C$ alone does not specify which lift preserves a cusp.
Also $\Pi_s$ is an extension, not automatically a split extension;
the construction has not supplied a section of $\Pi_s\to G_k$.
In particular the pullback of the diagonal cusp extension along a
nongeometric section is a virtual boundary datum, whose realizability
as a rational puncture has not been proved by writing the fibre product.

\begin{claim}{C.LIT3}{Proved: precise relation to Sun's virtual curves}{A14, C4}
The construction in C.LIT1 is already Sun's virtual fundamental group.
The two opened preprints do not assert that every virtual puncture is
the puncture at a rational point.
\end{claim}
\begin{proof}
Sun 2025 Definition~3.1 starts with a smooth surjective morphism over
a characteristic-zero field and a Galois section of the base, and defines
the virtual group by this pullback. Corollary~3.2 records its extension
and the independence of the geometric kernel from the section.
For a hyperbolic-curve family, Proposition~3.7 gives the fibre comparison,
and Corollary~3.8 identifies the virtual group with the actual fibre
when the section is geometric. This is a direct predecessor, rather
than a new construction of this paper.

Sun 2026 develops additional reconstruction and geometricity criteria.
Its introduction's Theorem~D corresponds to Corollary~5.8 in the body
(the introduction calls the reference a lemma). It concerns a number
field or a mixed-characteristic local field, virtual curves arising
from curves of types $(0,3)$ and $(1,3)$, and an open injection of their
virtual groups respecting their geometric kernels. Under those
hypotheses it forces the geometric index to be two and both sections
to be geometric. These restrictive hypotheses are not available for
an arbitrary section over a finitely generated extension of $\Q$.
\citep[Definition~3.1, Corollary~3.2, Proposition~3.7 and
Corollary~3.8]{LA_Sun2025}
\citep[Theorem~D and Corollary~5.8]{LA_Sun2026}
\end{proof}

\begin{claim}{C.LIT4}{Proved: the characteristic-zero recognition criterion}{C4}
Suppose $k$ is finitely generated over $\Q$, $C/k$ is hyperbolic,
and the base quotient retains the full arithmetic fundamental group,
or at least the point-separating Kummer maps used below. A cusp-preserving
isomorphism over the identity of this base quotient between the
punctures at $P,Q\in C(k)$ forces $P=Q$.

More generally, in the full base quotient, if the virtual puncture
$(\Pi_s,D_s)$ is isomorphic over $\Pi_C$, preserving the distinguished
cusp, to the actual puncture at $P$, then $s$ and $s_P$ are geometrically
conjugate. Conversely a geometric section has that puncture by C.LIT1.
\end{claim}
\begin{proof}
Projecting the cusp decomposition groups to $\Pi_C$ identifies the
point sections up to geometric conjugacy. If the smooth proper
compactification $\bar C$ has positive genus, abelianization gives
zero Kummer class for $[P-Q]\in J_{\bar C}(k)$. The
Mordell--Weil--Lang--N\'eron finite-generation theorem applies to
$J_{\bar C}(k)$. For a finitely generated
abelian group $A$, the map to all its finite-quotient completions is
injective: $\bigcap_{n\ge1}nA=0$. The Kummer exact sequences therefore
make the all-prime Kummer map injective on $J_{\bar C}(k)$. Hence $[P-Q]=0$.
A nonconstant function with divisor $P-Q$ would have degree one and
identify $\bar C$ with $\mathbf P^1$, impossible for positive genus; thus
$P=Q$.

If $\bar C$ has genus zero, extend $k$ finitely so that two boundary
points $a,b$ are rational. The existence of $P$ identifies $\bar C$
with $\mathbf P^1$. A degree-one coordinate with zero at $a$ and pole
at $b$ gives an injective map $C\to\mathbf G_m$ on rational points.
The full Kummer map on $k^\times$ is injective. To see this, let $K$
be the algebraic closure of $\Q$ in $k$, a number field. If an element
of $k^\times$ has an $n$th root in $k$ for every $n$, all its divisorial
valuations on a normal projective $K$-model are zero, so it lies in
$K^\times$. Its roots also lie in $K$, because they are algebraic over
$K$. Their valuations at every finite prime are zero; the finitely
generated unit group of $K$ has trivial intersection of all its
$n$th-power subgroups. Thus the element is one. Applying this to the
ratio of the coordinate values at $P,Q$ again proves $P=Q$.

For the virtual statement, the image of its distinguished cusp in
$\Pi_C$ is $s(G_k)$ by definition. An over-$\Pi_C$, cusp-preserving
comparison with the actual puncture makes this subgroup conjugate
to $s_P(G_k)$. Their projections to $G_k$ identify the corresponding
sections. Conversely C.LIT1 and C.LIT2 give the comparison when $s$
is geometric.
\end{proof}

This gives a recognition equivalence for the open characteristic-zero
question; it neither constructs a nongeometric section nor proves
that none exists. The finite-field proof of AG.40 used finiteness of
$J_C(k)$ to obtain a principal multiple of $P-Q$, constant-field point
counts, and the weight calculation eliminating nonconstant proper
Kummer classes. Those steps cannot simply be reused over $k$.
In a singleton-$l$ quotient, even the abelian Kummer map on $J_C(k)$
has the prime-to-$l$ torsion subgroup in its kernel. Thus exactness of
a chosen configuration quotient is not by itself a proof of point
separation for its base quotient. The full-group case, including
affine curves, is covered by the genus-zero and positive-genus
arguments above.\citep[\S1.2]{LA_CiperianiStix2015}
